%!TEX root = ../../main.tex
%% ===============================================================
%% 5.1 동기
%% 파일: chapters/05_label/01_motivation.tex
%% 상위: chapters/05_label/chapter.tex
%% 정의 라벨: sec:label-motivation
%% 참조 라벨: -
%% 포함 객체: -
%% 인용: -
%% 상태: 초안 (docs/OUTLINE.md 참고)
%% ===============================================================

\section{동기}
\label{sec:label-motivation}

가장 단순한 교전 결과 라벨은 라벨 창 안의 킬 수 차이 $\Delta_{\mathrm{kill}} = \mathrm{kills}_{\mathrm{blue}} - \mathrm{kills}_{\mathrm{red}}$의 부호이다. 그러나 킬 수는 교전의 \emph{가치}를 제대로 반영하지 못한다. 예를 들어 (i) 한 팀이 1 킬을 내주고 바론을 획득한 교전, (ii) 연승 중인 상대의 셧다운(shutdown) 킬 하나와 일반 킬 둘의 교환, (iii) 에이스(ace)로 끝난 교전은 킬 수만으로는 같은 값이거나 반대 부호가 된다. 프로 해설이 교전을 평가할 때 사용하는 ``이득'' 개념은 킬·오브젝트·골드가 혼합된 교환의 순가치이다. 본 논문은 이를 \emph{교환가치}(exchange value)로 형식화한다.
